\documentclass[
	a4paper,
	%DIV=13,
	smallheadings,
	%headinclude,
	%footinclude,
	%headsepline,
	%parindent,
	german,
	%captions=tableheading,
	%abstracton
	]
	{scrreprt}
\usepackage{blindtext}
\usepackage[T1]{fontenc}
\usepackage[utf8]{inputenc}
\usepackage[ngerman]{babel}
\usepackage[
	left	= 2 cm,
	right	= 2 cm,
	top		= 2 cm,
	bottom	= 4 cm
	]
	{geometry}
%\typearea[current]{calc}
\usepackage{lmodern}			
\usepackage{microtype}			
\usepackage[font=small,labelfont=bf]{caption}
\usepackage{amsmath,amssymb}
\usepackage{chemformula}
\usepackage{chemfig}
\usepackage{siunitx}
\DeclareSIUnit\year{a}
\usepackage{tabularx} 			
\usepackage{mdwlist}
\usepackage{caption}
\usepackage[babel,german=quotes]{csquotes}
\usepackage{xcolor}
\usepackage{hyperref}
\hypersetup{
		colorlinks	= true,
		linkcolor	= black,
		urlcolor	= blue,
		citecolor	= blue,
		pdftitle    = {Titel},
  		pdfsubject  = {CNC},
  		pdfauthor   = {Christopher Labisch},
  		pdfkeywords = {Tags} ,
  		pdfcreator  = {pdflatex},
  		pdfproducer = {LaTeX with hyperref}
	}
\usepackage[style=numeric-comp]{biblatex}
\def\bibfont{\footnotesize}
\bibliography{../Literatur/LiteraturCNC}

\title{CNC Eigenbau}
\subtitle{Open Source Projekt, Version 1.0}
\date{18.05.2017}
\author{Alexander Frosinn, \\ Dawid Konczak, \\ Christopher Labisch, \\ Alexander Lohberg}

\begin{document}


\maketitle
\tableofcontents
\chapter{Open Hardware CNC-Maschinen}
\begin{tabularx}{\textwidth}{XX}
\includegraphics[width=0.2\textwidth]{../Grafiken/Opensource.png} & \includegraphics[width=0.2\textwidth]{../Grafiken/Open-source-hardware-logo.png}
\end{tabularx}
\section{Generelles}
Eine kleine Bildergalerie zu industriellen und DIY-CNC-Maschinen stammt von Roman Black \autocite{:01}.
Er kommentiert einzelne Aspekte der Modelle wie Schlupf von Linearwagen, Antriebsarten, austauschbare Werkzeuge oder verwendete NEMA-Schrittmotoren.


\section{Arcitype}

Ein ausführliches CNC Projekt ist von DoES Liverpool beschrieben. \autocite{:Arcitype}


\section{HomoFaciens CNC}
Norbert Heinz (auch bekannt als HomoFaciens) hat einige Prototypen gefertigt.
Hier seien die neusten Entwicklungen angesprochen.
Als Designstudie wurde ein Plotter-Modell aus Pappe gebaut \autocite{:Heinz_CNC_3.0}.
Besonders die Steifigkeit und Geschwindigkeit als Schwachpunkte sind anschaulich erklärt.

In der Version 3.1 werden experimentell verschiedene Linearantriebe betrachtet \autocite{:Heinz_CNC_3.1}.
Diese werden mit einem alten Druckerantrieb, Bürstenmotoren oder Schrittmotoren angetrieben.
Als Sensoren zur Postionsbestimmung kommen Gabellichtschranken, eine optische Maus, oder ein Rotationsencoder zum Einsatz.
Die Ansteuerung erfolgt mit einem Arduino Uno.

In der Version 3.2 wird der Aspekt der Konstruktion näher erläutert \autocite{:Heinz_CNC_3.2}.
Besonders die Bearbeitung der Metallprofile wie Bohren, Biegen, Gewindeschneiden werden erläutert.
Die Spindel wird testweise über verschiedene Übersetzungs-Arten bzw. Zahnradgetriebe angetrieben.
Außerdem wird ein Kühlmittelkreislauf zur Verbesserung der Auflösung in der Tiefe (z-Achse) gezeigt. 
Vor der Inbetriebnahme für einen Test wird der Frästisch über Schrauben feinjustiert.
Nachteilig sind die große Anzahl an Verschraubungen und die mühselige Konstruktion an sich.
Der Preis für die Einzelteile wird mit ungefähr 350€ angegeben. 

\section{RepRap CNC}
Es gibt einige CNC-Versionen auf RepRap-Basis \autocite{:02}.

\section{MicroMill CNC}
Eine kleine CNC-Version im Schreibtischformat wurde kürzlich auf Kickstarter beworben \autocite{:kickstarter}.
Bei diesem Modell wird besonders Wert auf die Einfachheit und Flexibilität gelegt.
Ein besonderes Merkmal ist die Auswahl der Kugel- und Linearlager (Fa. Igus) und die Nutzung eines verfügbaren Wekrzeugs (Fa. Proxxon). 
Der Preis beläuft sich auf \$818, wobei nicht weiter aufgeschlüsselt wird. 

\chapter{Projektmanagement}

Die Herausforderung liegt in dem Projektmanagement, da es sich um ein recht umfangreiches Projekt handelt, das mit wenigen Ressourcen auskommen muss.
Viele Teile benötigt werden und eine lange zu erwartende Produktionsphase steht bevor.
Ein wesentliches Ziel ist die Projektphase so kurz wie möglich zu halten bei Gewährleistung einer guten Produktqualität.

Zu Beginn sind drei Personen involviert, die dieses Projekt unterstützen.
Für den 10.6.2017 ist ein erstes Treffen andauernd, Bei dem konzeptionellen Entwürfe erarbeitet und wichtige Rahmenbedingungen geklärt werden sollen.
Das Projekt als Open Source Hardware ausgelegt ist, soll es au Github dokumentiert werden.
Für meine persönliche Planung würde ich \texttt{Microsoft Project} benutzen, ich kann mir vorstellen, das außer mir niemand dieses Programm benutzt.
Daher wäre es gut, wenn es ein Programm oder eine Plattform gäbe, wo die Ressourcen geplant werden können.

Es könnte eine Art Fahrplan gestellt werden in Form eines Gantt-Diagramms mit Projektphasen und der Definition von Meilensteinen.
Zu Beginn steht die Erarbeitung von Konzepten im Vordergrund, wobei doch schon klar werden soll welche Anforderungen an die Maschine gestellt werden und wie man diese Anforderungen einfach erfüllen kann.
Dies betrifft insbesondere die Materialliste, die möglichst kurz ausfallen soll und die Verfügbarkeit von Materialien.
Ein weiterer wichtiger Aspekt ist die Durchführung der Arbeit an entsprechenden Orten mit entsprechendem Werkzeug.

\section{Workshops}
\subsection{1. CNC workshop -- 11.06.2017}
\paragraph{Brainstorming}
Zu Beginn haben Christopher und Dawid eine zweiminütige Präsentation der Idee gehalten.
Alexander L., Christopher, Dawid und Suleiman haben jeder für sich in 5 Minuten Anforderungen an die CNC auf Karten geschrieben.
Diese Anforderungen wurden vorgestellt, kategorisiert, priorisiert und konträre Anforderungen identifiziert.
Zuletzt hat Alexander L. die Ergebnisse in \texttt{Visual Understanding Environment} (\texttt{VUE}) übertragen.
Die \textit{Mindmap} \texttt{Fraese.vue} wurde im Verzeichnis \texttt{/Konzepte} abgelegt.

\subsection{2. CNC workshop -- 15.06.2017}
\paragraph{Bauraum}
Anhand dieser Mindmap wurden die Anforderungen an die Konstruktion konkretisiert.
Alexander L. hat die Erörterung moderiert.
Christopher und Dawid haben vorgestellt, welche Produkte mit der CNC-Fräse sie herstellen möchten.
Die Ergebnisse wurden am Whiteboard skizziert (\texttt{Bauraum- zu fräsende Objekte-Skizze.JPG} im Verzeichnis \texttt{/Konzepte}) und in die Mindmap schriftlich festgehalten.
Die Erörterung brachte hervor, dass eine Präzisionsfräse am besten die Anforderungen erfüllt.

Alexander F. hat angeregt, welche formalen Aspekte bei der Fräskopf-Auswahl hinsichtlich des zu bearbeitenden Materials zu beachten sind.
Zwei nützliche Hyperlinks zu einem ähnlichen Projekt\autocite{:Herbert} und Werkzeug-Angeboten\autocite{:ZuJeddeloh} wurden in die Mindmap  aufgenommen.

\paragraph{Fräskopf}
Es folgte eine Diskussion über mögliche Fräsköpfe.
Es wurden angreifende Kräfte (Normalkraft, Querkraft, ...) und Materialeigenschaften der Werkstücke (Carbon, Aluminium, Kupfer, Silber, Messing, Maschinenmessing, ...)  besprochen.
Es folgen weitere Recherchen über die zu erwartenden Kräfte (Dawid) und Software-Ansteuerung (Christopher).

Das weitere Projekt-Vorhaben wurde mit \texttt{ProjectLibre} als \textit{Gantt-Diagramm} (\texttt{Fräsenbau20170615.pod}) im Verzeichnis \texttt{/Konzepte} festgehalten. 


\section{Projektphasen}
\subsection{Definition von Projektphasen und Meilensteinen}
\subsection{Anforderungsprofil}
\subsection{Recherche}
\subsection{Konzeptionierung}
\subsection{Preisvergleich}

\chapter{Komponenten}
\section{Koordinatentisch}
	\subsection{Aluminiumsystemprofile}
	\subsection{Tisch}
	\subsection{Werkstück-Befestigung}
	\subsection{Kühlmittel-Auffangbecken}
\section{Antrieb}
	\subsection{Antriebswelle}
	\subsection{Antriebsmotoren}
		\section{Bürstenmotoren}
		\section{Schrittmotoren}
\section{Werkzeug}
	\subsection{Werkzeugantrieb}
	\subsection{Bohrmaschine}
		\subsection{Bohrfutter}
		\subsection{Bohrer}
\section{Gerüst}
	\subsection{L-Form}
	\subsection{C-Form}
\subsection{Stahlprofile}
\subsection{Schutz}
\section{Kühlkreislauf}
	\subsection{Kühlmittel}
	\subsection{Schläuche}
	\subsection{Pumpe}
	\subsection{Filter}
\section{Elektronik-Hardware}
	\subsection{Netzteil}
	\subsection{Kabel}
	\subsection{Kabelführung}
	\subsection{Display}
	\subsection{PC-Schnittstelle}
	\subsection{Microcontroller}
	\subsection{Treiber}
	\subsection{Relais}
\section{Elektronik-Software}
Die Schrittmotoren und die Fräse sollen über einen Arduino angesteuert werden sollen.
Dafür ist es nötig, dass ein Programm ein Modell des Werkstücks in Pfade umwandelt.
Die Pfade sollen parametrisiert abgefahren werden können.

\subsection{Grbl -- open-source CNC controller}
Grbl ist eine open-source Software zur Steuerung von Maschinenbewegungen mit Hilfe eines Arduino Uno \autocite{:Sonny2017}.

Es gibt verschiedne Hobby- und professionelle Projekte, die Grbl als Software benutzen \autocite{:Herbert,:kickstarter,:Meissner}.

Viele open-source 3D-Drucker nutzen diese software, außerdem gibt es Modifikationen für ähnliche Anwendungen.
Die Software verfügt über ein grafisches User-Interface (GUI), ermöglicht aber auch die Eingabe über eine Konsole.
Grbl verarbeitet \textit{g-Code}.

Grbl kann kartesische Koordinaten abfahren (x, y, z). 
Mit Grbl ist es jedoch (noch) nicht möglich eine Rotationsbewegung (d.h. in Polarkoordinaten) umzusetzen.

\chapter{Methoden}

\section{Modellieren}
	\subsection{FreeCAD}
	\subsection{Illustrator}
	\subsection{Excel}
\section{Konstruktion}
Bei der Konstruktion ist zu beachten, dass die Module nicht zu großen Drehmomenten ausgesetzt sind, da sonst Spiel in die Mechanik kommt.
Deshalb ist es wichtig Torsionskräfte und Hebelkräfte zu kompensieren.
Dazu eignen sich z.B. Fachwerke oder Knotenbleche.
	\subsection{Fachwerk}
	\subsection{Knotenbleche}
Ein alter Dachdecker-Spruch besagt: \quote{Viereck vergeht, Dreieck besteht.}
	\subsection{Stecken}
	\subsection{Langloch-Verbinden}
	\subsection{Schrauben}
		\subsection{Schraubverbindung}
	\subsection{Lager}
		\subsection{Linearlager}
		\subsection{Axiallager}
\section{Bohren}
\section{Schweißen}
	\section{WIG}
	\section{MAG}
\section{Schneiden}
\subsection{Gewindeschneiden}


\printbibliography


\end{document} 